% ============================================================
% Section 3: Methodology
% ============================================================

\section{Methodology}
\label{sec:method}

\subsection{Dataset}
\label{subsec:dataset}

We train and evaluate DocVerify exclusively on \textbf{FantasyID}~\cite{fantasyid2024},
a publicly released benchmark for identity document forgery detection and
localization developed by Idiap Research Institute for the DeepID 2025
Challenge~\cite{deepid2025challenge}. FantasyID comprises approximately 3\,284
labeled images of synthetic fantasy identity documents in two classes: \textit{bonafide}
(genuine, unmodified) and \textit{attack} (tampered). Attack images are produced
via two representative manipulation types: (i)~\textbf{face swap}, replacing the
portrait photograph with a different identity, and (ii)~\textbf{text inpainting},
overwriting text fields (name, date of birth, document number) with forged content.

Each image is accompanied by a JSON annotation file specifying bounding-box
regions with a \texttt{region\allowbreak\_provenance:\allowbreak~altered} field that identifies the
tampered areas. We convert these annotations to binary segmentation masks by
rasterizing the altered bounding boxes at the training resolution
($224 \times 224$ pixels). Bonafide images have all-zero masks. We combine the FantasyID \texttt{train}
and \texttt{test} partitions into a single pool of 3\,284 images and apply a
\texttt{GroupShuffleSplit} (grouping by document identity, $\text{test\_size}=0.15$,
\texttt{random\_state}=42) to obtain approximately 2\,791 development images and
493 blind holdout images, ensuring that all captures of the same document identity
remain in the same partition and preventing any form of data leakage between splits. The holdout set is accessed only for the final blind
evaluation (Section~\ref{sec:blindtest}).

\subsection{Architecture}
\label{subsec:arch}

DocVerify is a multi-task encoder-decoder network designed for simultaneous
document-level classification and pixel-level forgery localization.
Figure~\ref{fig:architecture} shows the overall structure.

\begin{figure*}[!t]
  \centering
  \includegraphics[width=\textwidth]{fig0_architecture_compact}
  \caption{DocVerify architecture. The encoder maps a
    $3{\times}224{\times}224$ input to a $256{\times}7{\times}7$ bottleneck through
    six convolutional blocks ($8\!\to\!256$ channels). The U-Net decoder reconstructs
    the forgery mask via skip connections; the classification head outputs a
    document-level tampered/genuine score.}
  \label{fig:architecture}
\end{figure*}

\subsubsection{Encoder}

The encoder is a lightweight six-block convolutional backbone inspired by the
dense CNN (D-CNN) architecture proposed by Patel et al.~\cite{patel2023dcnn}
for deepfake image detection, adapting its increasing channel-width progression
and mixed pooling strategy to the document-forensics domain.
It employs an increasing channel progression of $8 \to 16 \to 32 \to 64 \to 128 \to 256$
filters. All convolutional layers use LeakyReLU activations (negative
slope~$\alpha=0.2$)~\cite{leakyrelu2015} and batch normalization~\cite{batchnorm2015}.
Blocks 1--4 use $3 \times 3$ convolutions; blocks 5 and~6 use $5 \times 5$
convolutions to capture wider spatial context at lower feature-map resolutions.
Spatial downsampling is applied after each block: average pooling after blocks 2
and~3 (to preserve boundary information for localization), and max pooling after
blocks 4 and~5 (to achieve stronger translation invariance for classification).
The encoder maps a $3 \times 224 \times 224$ input to a bottleneck of $256 \times 7 \times 7$.

Skip connections are extracted from each block prior to pooling, yielding feature
maps at five resolutions: $224\times224$ (16~ch), $112\times112$ (32~ch),
$56\times56$ (64~ch), $28\times28$ (128~ch), and $14\times14$ (256~ch).

\subsubsection{Classification Head}

The classification head applies global average pooling to the $7 \times 7$
bottleneck, yielding a 256-dimensional feature vector. This is passed through
four fully-connected layers (256\,$\to$\,32\,$\to$\,16\,$\to$\,16\,$\to$\,1)
interleaved with dropout and LeakyReLU activations, producing a single
classification logit $\hat{y}_{\text{cls}}$.

\subsubsection{Segmentation Head (U-Net Decoder)}

The segmentation head reconstructs the spatial mask through a five-stage U-Net
decoder~\cite{unet2015}. A $1 \times 1$ projection convolution first maps the
bottleneck to $C_{\text{dec}}$ channels (a tunable hyperparameter, see
Section~\ref{subsec:hpo}). Each decoder block bilinearly upsamples by a factor of
two, concatenates the corresponding encoder skip connection, and applies two
$3 \times 3$ convolutions with LeakyReLU activations—mirroring the standard
U-Net design. The channel width halves at each stage from $C_{\text{dec}}$ down
to $C_{\text{dec}}/16$. A final $1 \times 1$ convolution produces the segmentation
logit $\hat{Y}_{\text{mask}}$ at the original input resolution.

Both output logits are kept raw (no sigmoid at inference time); sigmoid is applied
numerically within the loss functions for numerical stability.

\subsection{Loss Functions}
\label{subsec:loss}

The combined training loss is:
\begin{equation}
  \mathcal{L} = \mathcal{L}_{\text{cls}} + \lambda_{\text{mask}} \,
                \mathcal{L}_{\text{seg}},
  \label{eq:total_loss}
\end{equation}
where $\lambda_{\text{mask}} \in [0.5, 3.0]$ is a tunable weight included in the
HPO search space (Section~\ref{subsec:hpo}).

\paragraph{Classification loss.}
$\mathcal{L}_{\text{cls}} = \text{BCEWithLogits}(\hat{y}_{\text{cls}},\, y_{\text{cls}})$,
where $y_{\text{cls}} \in \{0, 1\}$ is the document-level label.

\paragraph{Segmentation loss.}
$\mathcal{L}_{\text{seg}} = \text{BCEWithLogits}(\hat{Y}_{\text{mask}},\, Y_{\text{mask}})
  + \mathcal{L}_{\text{Dice}}(\sigma(\hat{Y}_{\text{mask}}),\, Y_{\text{mask}})$,
where $\sigma(\cdot)$ denotes the sigmoid function and the Dice
loss~\cite{dice1945} is computed per image and averaged over the batch:
\begin{equation}
  \mathcal{L}_{\text{Dice}} = 1 - \frac{2 \sum_i p_i g_i + \epsilon}
                                        {\sum_i p_i + \sum_i g_i + \epsilon},
  \quad \epsilon = 10^{-6}.
\end{equation}
The BCE--Dice combination was introduced for volumetric medical image
segmentation~\cite{milletari2016vnet} and has since become standard practice for
pixel-level tasks with severe class imbalance~\cite{unet2015}; it directly addresses
the foreground/background imbalance caused by small forgery regions such as a single
altered text field.

\paragraph{Training details.}
All models are trained with AdamW~\cite{adamw2019} and ReduceLROnPlateau (factor
$0.5$, patience~3). Gradient clipping with a maximum $\ell_2$ norm of~1.0 is
applied at every step. Mixed-precision training (FP16 forward pass, FP32 gradient
accumulation via the PyTorch AMP GradScaler~\cite{micikevicius2018mixed,pytorch2019})
is used throughout, reducing memory usage without affecting convergence.
Input images are normalized to $[0, 1]$ and resized to $224 \times 224$ pixels.

\subsection{Multi-Objective HPO}
\label{subsec:hpo}

Our HPO formulation treats detection quality and localization quality as two
independent objectives to be simultaneously maximized:
\begin{equation}
  \max_{\boldsymbol{\theta} \in \Theta}\; \bigl(\text{PR-AUC}(\boldsymbol{\theta}),\;
  \text{Dice}(\boldsymbol{\theta})\bigr),
  \label{eq:moopt}
\end{equation}
where $\boldsymbol{\theta}$ is a hyperparameter configuration and $\Theta$ is the
search space defined in Table~\ref{tab:hpo_space}.

\begin{table*}[!t]
  \centering
  \caption{HPO search space}
  \label{tab:hpo_space}
  \begin{tabular}{llcc}
    \toprule
    Parameter & Distribution & Range & Scale \\
    \midrule
    $\eta$ (learning rate) & log-uniform & $[5\times10^{-5},\; 9\times10^{-4}]$ & log \\
    $\lambda_{\text{wd}}$ (weight decay) & log-uniform & $[10^{-7},\; 10^{-4}]$ & log \\
    $p$ (dropout rate) & uniform & $[0.1,\; 0.4]$ & linear \\
    $C_{\text{dec}}$ (decoder channels) & categorical & $\{96, 128, 192, 256\}$ & — \\
    $\lambda_{\text{mask}}$ (mask loss weight) & uniform & $[0.5,\; 3.0]$ & linear \\
    \bottomrule
  \end{tabular}
\end{table*}

The LeakyReLU slope $\alpha$ was initially included in the search space but removed
after a preliminary study of 250 trials showed no signal on validation performance
across the range $[0.1, 0.4]$; it is fixed at $\alpha = 0.2$ throughout.

\subsubsection{MOTPE Sampler}

We use the Multiobjective Tree-structured Parzen Estimator
(MOTPE)~\cite{motpe2022} as implemented in Optuna~\cite{optuna2019} with NSGA-II~\cite{nsga2_2002}
as a fallback sampler. MOTPE extends the single-objective TPE by modeling
the joint distribution of hyperparameter configurations in the objective space,
efficiently focusing the search toward the Pareto-optimal frontier without
requiring an explicit scalarization function. This is key for our setting, where
neither objective should be privileged a priori.

\subsubsection{Pareto-Based Model Selection}

After running 50 Optuna trials on the inner validation split of each outer fold,
we identify the set of Pareto-optimal trials (non-dominated solutions). From this
set, we select the configuration $\boldsymbol{\theta}^*$ that minimizes the
Euclidean distance to the ideal point $(1, 1)$ in the objective space:
\begin{equation}
  \boldsymbol{\theta}^* = \argmin_{\boldsymbol{\theta} \in \mathcal{P}}
    \sqrt{(1 - \text{PR-AUC})^2 + (1 - \text{Dice})^2},
  \label{eq:ideal_point}
\end{equation}
where $\mathcal{P}$ denotes the Pareto front. This criterion requires no
user-specified weighting and is consistent with the DeepID 2025 Challenge
evaluation protocol, which ranks systems independently on Track~1 (detection)
and Track~2 (localization), assigning equal importance to both.

Figure~\ref{fig:pareto} illustrates the objective space across all 500 HPO trials
(10 folds $\times$ 50 trials), with the Pareto-optimal trials highlighted.

\begin{figure*}[!t]
  \centering
  \includegraphics[width=\textwidth]{fig1_pareto_front}
  \caption{Validation objective space across all 500 HPO trials (10 outer folds
    $\times$ 50 Optuna MOTPE trials per fold). Grey points: evaluated configurations;
    blue points: Pareto-optimal (non-dominated) trials; dashed line: global Pareto
    frontier. The red diamond at $(1,~1)$ marks the ideal point; the annotated
    point is the selected configuration (minimum Euclidean distance $d = 0.316$).}
  \label{fig:pareto}
\end{figure*}

\subsection{Nested Cross-Validation}
\label{subsec:ncv}

To obtain unbiased performance estimates while keeping HPO blind to the test data,
we adopt a three-level nested cross-validation (NCV) scheme illustrated in
Figure~\ref{fig:nested_cv}.

\paragraph{Outer loop (10 folds).}
The FantasyID training set is partitioned into 10 stratified folds. Each fold
serves once as the test set for a held-out performance estimate.

\paragraph{Inner loop (5 folds, HPO).}
For each outer fold, the remaining 9/10 of the data is further split into 5 inner
folds. The 50 Optuna MOTPE trials are run on these inner splits, with each trial
training for up to $E_{\text{trial}} = 15$ epochs with patience~$= 1$ on the inner
validation loss. The optimal configuration $\boldsymbol{\theta}^*$ is selected via
(\ref{eq:ideal_point}).

\paragraph{Final model training.}
The selected $\boldsymbol{\theta}^*$ is used to retrain the model from scratch on
all 9/10 outer-fold training data for up to $E_{\text{final}} = 100$ epochs with
early stopping (patience~$= 12$, monitoring the distance-to-ideal criterion on a
held-out inner validation slice). The trained model is then evaluated on the
held-out outer-fold test set to produce one unbiased performance estimate.

This design ensures that neither the HPO process nor the final training sees the
test fold, eliminating the optimistic bias that arises when hyperparameter selection
and performance estimation share data~\cite{varma2006bias}. The total number of
model trainings is $10 \times 5 \times 50 = 2\,500$ (HPO) plus $10$ final models.
